\section{Hugging Face Model Card}
\label{app:hf_model_card}

We release the ladder on Hugging Face.\footnote{Link removed for anonymity.} This section includes the model card that will accompany the models, following the Hugging Face model card template \citep{ozoani_model_2022}.

% https://huggingface.co/docs/hub/model-card-annotated
% Template: https://github.com/huggingface/huggingface_hub/blob/main/src/huggingface_hub/templates/modelcard_template.md

\newpage

\subsection*{Model Card for the Multilingual Signal-and-Noise Model Ladder}

A ladder of \nruns decoder-only language models (90M to 3B non-embedding parameters) pretrained from scratch on controlled multilingual data mixtures, released with all intermediate checkpoints to study which evaluation benchmarks give reliable signal at small scale.

\subsection*{Model Details}

\subsubsection*{Model Description}

The ladder varies one design decision at a time around a fixed recipe: model size $N \in \{$90M, 175M, 350M, 600M, 1B, 1.7B, 3B$\}$, the number of languages in the mixture $K \in \{1, 2, 8, 15, 30, 50\}$, the language list (data scheme A, B or C), the sampling temperature ($T = 1$ or 3), the depth-to-width ratio (deep or shallow) and the MLP activation (xIELU or SwiGLU). Every model trains on $D = 100N$ tokens (5$\times$ Chinchilla) with half of the tokens in English. Each run is one model repository; its checkpoints are branches.

\begin{itemize}
    \item \textbf{Developed by:} Anonymous (removed for review)
    \item \textbf{Funded by:} Anonymous (removed for review)
    \item \textbf{Shared by:} Anonymous (removed for review)
    \item \textbf{Model type:} decoder-only transformer, causal language model, pretrained only (no instruction tuning)
    \item \textbf{Language(s) (NLP):} English plus up to 49 other languages (Appendix~\ref{app:data-mixtures}, Table~\ref{tab:app-languages}); each model's card lists its own $K$ languages
    \item \textbf{License:} \checknum{Apache 2.0}
\end{itemize}

\subsubsection*{Model Sources}

\begin{itemize}
    \item \textbf{Paper:} Link removed for anonymity.
    \item \textbf{Repository:} Link removed for anonymity.
    \item \textbf{Website:} Link removed for anonymity.
\end{itemize}

\subsection*{Uses}

\subsubsection*{Direct Use}

Research on pretraining and evaluation: ablations of data mixtures and architectures at small scale, scaling-law fits across size and training, checkpoint-level analyses of benchmark signal and noise, and as reference proxies when choosing benchmarks for multilingual pretraining runs.

\subsubsection*{Downstream Use}

Starting points for continued pretraining or fine-tuning experiments that need many comparable models trained on known data. The models are not adapted to follow instructions.

\subsubsection*{Out-of-Scope Use}

Deployment in any user-facing application, generation of factual content, or use in high-stakes settings. The models are small, untuned and not safety-aligned; their outputs are unreliable and can be harmful.

\subsection*{Bias, Risks, and Limitations}

The models inherit the biases of web text (DCLM and FineWeb-2) and of its language distribution: English is half of every mixture, and the other languages follow the per-language token shares of the mixture, so lower-resource languages are seen far less. Many multilingual benchmark tasks are at chance for these model sizes (Section~\ref{subsec:rq0}). The 90M and 175M rungs train at a reduced batch (Appendix~\ref{app:model-configurations}).

\subsubsection*{Recommendations}

Compare models only within the ladder's controlled settings, read results with the above-chance gate and the seed replicates as the noise floor, and do not treat any checkpoint as a usable assistant.

\subsection*{How to Get Started with the Model}

\begin{small}
\begin{verbatim}
from transformers import AutoModelForCausalLM, AutoTokenizer

repo = "<org>/lm-1.7B-L8-deep-seed1904"   # link removed for anonymity
tok = AutoTokenizer.from_pretrained(repo)
model = AutoModelForCausalLM.from_pretrained(repo, revision="main")
# intermediate checkpoints: revision="iter_<step>"
\end{verbatim}
\end{small}

\subsection*{Training Details}

\subsubsection*{Training Data}

English: DCLM-Edu; the two alternative $K = 1$ builds use DCLM without the educational filter and FineWeb up to 2022. Other languages: FineWeb-2, with the language lists and token shares of Appendix~\ref{app:data-mixtures}. Each mixture holds the English share at 50\% and varies only the other languages.

\subsubsection*{Training Procedure}

\paragraph{Preprocessing}

Tokenized with the 131,072-token Apertus tokenizer; documents are packed into 4,096-token sequences. Validation documents used for bits-per-byte are held out of every training mixture.

\paragraph{Training Hyperparameters}

\begin{itemize}
    \item \textbf{Training regime:} bf16 mixed precision
    \item \textbf{Optimizer:} AdEMAMix ($\beta_1 = 0.9$, $\beta_2 = 0.999$, $\beta_3 = 0.9999$, $\alpha = 8$), weight decay 0.1, gradient clipping 0.1
    \item \textbf{Schedule:} warmup-stable-decay, with a 1$-\sqrt{\cdot}$ cooldown over the last 20\% of training
    \item \textbf{Batch:} 504 sequences of 4,096 tokens from 350M up; 84 at 90M and 168 at 175M
    \item \textbf{Learning rate, depth and width per size:} Table~\ref{tab:ladder}
\end{itemize}

\paragraph{Speeds, Sizes, Times}

20 checkpoints per run (40 at 1B, 60 at 1.7B and 3B). Pretraining the ladder used \npretrainnodehours node-hours.

\subsection*{Evaluation}

\subsubsection*{Testing Data, Factors \& Metrics}

\paragraph{Testing Data}

\nbenchmarks multilingual benchmark families plus English benchmarks in the LM Evaluation Harness (\nbenchmarktasks benchmark--language tasks), and held-out bits-per-byte on 100 languages (Appendix~\ref{app:evaluation-protocol}).

\paragraph{Factors}

Model size, number of training languages, data scheme, temperature, depth, activation, seed and training checkpoint.

\paragraph{Metrics}

Accuracy (or normalized accuracy) per benchmark--language task, read against its chance level with a one-sided Wilson bound; bits-per-byte per language; decision accuracy and signal-to-noise ratio across the ladder.

\subsubsection*{Results}

\paragraph{Summary}

Section~\ref{sec:results} and the per-analysis appendices report the results; every number is regenerated from the released evaluation logs.

\subsection*{Environmental Impact}

\begin{itemize}
    \item \textbf{Hardware Type:} NVIDIA GH200 (4 per node)
    \item \textbf{Hours used:} \ngpuhours GPU-hours (training and evaluation)
    \item \textbf{Cloud Provider:} on-premise HPC (removed for review)
    \item \textbf{Compute Region:} removed for review
    \item \textbf{Carbon Emitted:} \checknum{not estimated}
\end{itemize}

\subsection*{Technical Specifications}

\subsubsection*{Model Architecture and Objective}

Decoder-only transformer trained with next-token prediction: grouped-query attention (4 query heads per key-value head, head dimension 64), rotary position embeddings, RMSNorm, QK-norm, xIELU (or SwiGLU) MLP, no biases, tied input and output embeddings, 4,096-token context.

\subsubsection*{Compute Infrastructure}

\paragraph{Hardware}

GH200 nodes with four GPUs each.

\paragraph{Software}

Megatron-LM (training), Hugging Face Transformers (released checkpoints), LM Evaluation Harness (evaluation).

\subsection*{Citation}

\textbf{BibTeX:} removed for anonymity.

\subsection*{Model Card Authors}

Anonymous (removed for review).

\subsection*{Model Card Contact}

Removed for anonymity.
